%% ---------------------------------------------------------------------------
%% wildtag.ai - Methods in Ecology and Evolution, Practical Tools
%% Skeleton draft. Compile with pdfLaTeX + Biber.
%%
%% Bibliography: BibTeX (not biber), using the BES's own besjournals style.
%% On Overleaf the default pdfLaTeX + BibTeX setting handles this automatically.
%%
%% Style options:
%%   \usepackage[anonymous]{mee}  hide authors for double-anonymised review
%%   \usepackage[final]{mee}      strip line numbers, notes and (cite) markers
%%
%% Colour key in the draft PDF:
%%   orange italics  = \note{}  guidance for the authors, delete before submitting
%%   green (cite)    = \cw      a citation is needed here
%%   red [BRACKETS]  = \tbd{}   a number or name still to be decided
%% ---------------------------------------------------------------------------
\documentclass[11pt]{article}
\usepackage{mee}

\meearticletype{Practical Tools}
\meetitle{wildtag.ai: a fully offline desktop workflow linking automated
species identification, distributed volunteer validation and
standards-compliant output for camera-trap surveys}
\meerunning{wildtag.ai: an offline camera-trap workflow}
\meeauthors{\tbd{Author One}$^{1}$, \tbd{Author Two}$^{2}$, \tbd{Author Three}$^{1}$}
\meeaffils{$^{1}$\tbd{Affiliation}. $^{2}$\tbd{Affiliation}.}
\meecorresponding{\tbd{email@institution.ac.uk}}

\begin{document}

\makemeetitle

\note{References use the British Ecological Society's own BibTeX style
(besjournals), which ships with TeX Live and needs no extra files. Cite with
natbib commands: \textbackslash citep\{key\} gives (Author, 2024) and
\textbackslash citet\{key\} gives Author (2024). Each green (cite) marker below
marks a place where such a citation is needed.}

\note{Alternative shorter title if the editors want one: ``wildtag.ai: from
camera-trap images to a validated, standards-compliant dataset, entirely
offline''.}

\note{Article-type flag. MEE also runs Applications (software and packages) and
Workflows (assemblages of existing software into pipelines for large datasets).
wildtag is software that assembles existing models, so an editor may reclassify
it. All three types share a 3,000-4,000 word limit and this structure, so the
cost of reclassification is low, but decide deliberately which you target.}

%% ---------------------------------------------------------------------------
\begin{meeabstract}

\item Camera traps now generate datasets far larger than research teams can
annotate by hand \cw. Automated detection and classification models have become
accurate enough for routine use \cw, but converting their raw predictions into a
verified, shareable dataset still requires scripting, cloud accounts or manual
file handling that non-specialist field teams cannot readily perform \cw.
\note{The gap is the ``last mile'' between a model prediction and a trustworthy,
standards-compliant dataset. Say this plainly; it is the paper's reason to exist.}

\item We present \texttt{wildtag.ai}, a free desktop application that joins the
full sequence into a single offline workflow: automated detection and species
classification, structured in-application validation, a self-contained package
for distributing images to volunteers for verification, and export to the
Camtrap Data Package (Camtrap DP) standard \cw. Models are bundled and hosted
locally through a small registry, so no coding, cloud service or internet
connection is required for the European workflow.

\item We describe the architecture and demonstrate the end-to-end pipeline on a
reproducible demonstration deployment of 10 sites and \tbd{1,000} images.
\tbd{Insert one headline quantitative result from a real run: the proportion of
frames cleared automatically as empty, reviewer time saved, or volunteer
throughput.}

\item \texttt{wildtag.ai} lowers the technical barrier to closing the loop
between automated prediction and a verified, standards-compliant camera-trap
dataset. Because the model registry and the Camtrap DP output are model- and
region-agnostic, the same workflow extends beyond European mammals to any
detector-classifier pair and any regional survey.
\note{This final sentence carries the ``adapted for more general use''
requirement that MEE states explicitly for Practical Tools. Do not cut it.}

\end{meeabstract}

\meekeywords{camera trap; species identification; deep learning; citizen
science; data validation; Camtrap DP; biodiversity monitoring; offline software}

%% ---------------------------------------------------------------------------
\section*{Data and code for peer review}

The \texttt{wildtag.ai} source code, the model registry, the demonstration-project
generator and the generated fixture are available for review at \tbd{private
repository or Zenodo DOI}. The application is distributed at \tbd{installer URL}
and documented in a user manual at \tbd{manual URL}. Model weights are
redistributed under their respective licences \cw.
\note{Confirm you are permitted to redistribute the DeepFaune weights
(CC-BY-NC-SA-4.0) inside the bundle. If that is uncertain, change the design so
wildtag fetches them on first use, and say so here. A reviewer may check this.}

%% ---------------------------------------------------------------------------
\section{Introduction}
\note{Target 700-800 words. Three beats. This section decides whether a reviewer
believes the paper is novel, so the third beat must land hard.}

\subsection*{Beat 1: automated prediction is largely a solved problem}
Camera-trap studies routinely collect hundreds of thousands of images, the
majority of which contain no animal \cw. General-purpose detectors \cw and
regional or global species classifiers \cw now perform well enough to be
embedded in routine practice \cw.
\note{Establish clearly that automated prediction is not what this paper claims
to advance. Conceding this early buys credibility for the claim that follows.}

\subsection*{Beat 2: the unsolved ``last mile''}
A prediction is not a dataset. Between a model's output and a publishable,
reusable dataset lie three tasks that remain awkward for teams without
programming skills: reviewing and correcting predictions at scale; sharing the
verification workload, for example with citizen-science volunteers \cw, without
requiring them to install a software stack; and exporting to a community
standard so that the result is interoperable \cw. Existing tools address parts
of this. Cloud platforms require accounts, uploads and connectivity \cw. Offline
desktop identifiers run models without coding but stop at prediction and export,
leaving validation and volunteer distribution to the user \cw. Dedicated review
tools support manual annotation well but are not coupled to modern classifiers
or to a distributed-validation loop \cw.
\note{This paragraph is the crux of the paper. Cite each comparator fairly and
state precisely what it does and does not do; Table 1 then does the heavy
lifting. Being scrupulously fair to AddaxAI/EcoAssist here is essential, because
it is the closest tool and the first thing a reviewer will raise.}

\subsection*{Beat 3: what wildtag.ai adds}
The contribution of \texttt{wildtag.ai} is not a new model but an integration.
To our knowledge it is the first free desktop tool that carries a survey from
raw images to a validated, standards-compliant dataset entirely offline,
including a built-in mechanism for parcelling verification work out to
volunteers who need install nothing. The design rests on four pillars: local
model hosting through a simple registry; structured in-application validation;
self-contained distributed volunteer validation; and native export to a
community data standard \cw. Because the first and fourth pillars are model- and
region-agnostic, the workflow is a template rather than a single-region tool.
\note{Verify the ``first'' claim against the current releases of every comparator
before submission. If any has added a comparable feature, soften to ``one of
few'' and lean on the completeness of the end-to-end offline integration
instead. An overstated priority claim is the easiest way to lose a reviewer.}

%% ---------------------------------------------------------------------------
\section{Materials and Methods}
\note{Target 1,200-1,400 words. This is the heart of a Practical Tools paper:
the reader should finish it able to reproduce or adapt the approach.}

\subsection{Design goals and audience}
\texttt{wildtag.ai} was built under four constraints drawn from the working
conditions of its intended users: no programming background; commodity Windows
laptops rather than servers or clusters; datasets of $10^{5}$ images or more;
and unreliable or administratively restricted internet access \cw.
\note{These constraints justify every subsequent design decision (offline,
no-install, checkpointed) and pre-empt the obvious reviewer question, ``why not
simply use a cloud platform?''.}

\subsection{Overall architecture}
The application separates a lightweight graphical interface, which depends only
on the Python standard library, from a bundled inference environment invoked as
a subprocess. The interface therefore never imports a deep-learning framework,
which keeps it responsive and allows a minimal, 28\,MB runtime to be shipped to
volunteers who only need the validation pane. Figure~\ref{fig:arch} shows the
architecture and the flow of data. Projects follow an opinionated folder
structure (site-wise image folders, a deployment table, and generated validation,
export and run-summary outputs), which makes a project self-documenting and
allows every subsequent step to locate its inputs without configuration.

\subsection{Model hosting and the registry}
Models are described in a plain-data registry recording an identifier,
architecture, weights, licence, citation and dependencies. Each classifier
supplies its own detector, so the pipeline has no hard dependency on any single
detection model \cw. Two back ends are currently provided: a European classifier
paired with a bundled lightweight detector, which runs entirely offline \cw, and
a global classifier executed as a self-contained subprocess whose weights are
cached locally on first use \cw. Adding a further model requires a registry entry
and an inference module implementing two functions.
\note{State the extensibility explicitly and concretely. This is the mechanism
by which the tool ``can be adapted for more general use'', which the article type
demands, so make it impossible to miss.}

\subsection{Running at field scale}
Several features exist specifically to make very large runs survivable on modest
hardware, and they are as much a part of the tool's practical contribution as the
models it wraps. Progress is streamed from the inference subprocess so that a
multi-day run reports continuously. Results are checkpointed periodically, so a
crash or a power failure resumes rather than restarts. A stop request finishes
the current batch and saves, behaving like a pause rather than a kill. The sort
of images into species folders is non-destructive and resumable, so an
interrupted sort is completed rather than repeated and no completed work is
deleted. Long file paths, which arise routinely from nested site and camera
folders on Windows, are handled explicitly.
\note{Motivate this section with the real large deployment (approximately 473,000
images across 103 sites) that drove each of these features. Reviewers of a
practical-tools paper reward exactly this kind of hard-won field engineering, so
do not bury it.}

\subsection{In-application validation}
The validation pane presents predictions as a tiled gallery, one species at a
time, with bounding boxes drawn and sibling detections in the same frame shown
for context. Corrections are made against a controlled label vocabulary, singly
or in multi-selected batches, and are written back to per-species tables so that
validation produces data rather than merely rearranging files. Human detections
are pixelated by default in accordance with data-protection requirements \cw.
Figure~\ref{fig:ui}a shows the interface.

\subsection{Distributed volunteer validation}
\note{This is the flagship section and the paper's signature contribution.
Describe it in enough detail that a reader could adopt the pattern in another
tool.}
Verification effort is the principal bottleneck once prediction is automated \cw.
\texttt{wildtag.ai} therefore packages predicted-species images into
self-contained archives of configurable size, each bundling the application and a
minimal runtime, so that a volunteer extracts one file, runs it with a single
action, validates images with no installation, no administrative rights and no
internet connection, and returns the archive. A collection step then merges the
returned validations back into the project and tracks which batches remain
outstanding. Figure~\ref{fig:loop} illustrates the loop. Archive size is
configurable because the two realistic modes of use pull in opposite directions:
small archives can be emailed or hosted for remote volunteers, whereas larger
folders suit paid annotators working from local storage.

\subsection{Standards-compliant export}
Completed projects are exported as a Camtrap Data Package, comprising deployment,
media and observation tables with an accompanying metadata descriptor \cw. The
output is therefore interoperable with the wider camera-trap and biodiversity
data ecosystem, including aggregators \cw, without further scripting.

\subsection{Demonstration deployment}
To provide a reproducible worked example, we generated a demonstration project
of 10 sites and approximately 100 images per site, with an accompanying
deployment table and a schema-complete results table.
\note{Be explicit and honest here: the demonstration images are placeholder
frames, so the species counts they yield are synthetic and illustrative. They
demonstrate that the plumbing works and that the example is reproducible; they do
not demonstrate a result. Before submission, run wildtag on a real dataset and
regenerate Table 2 and Figure 3b from genuine model output. A reviewer who
discovers synthetic counts presented as findings will reject the paper, so this
substitution is not optional.}

%% ---------------------------------------------------------------------------
\section{Results}
\note{Target 600-700 words. Report the demonstration as a walkthrough of the four
pillars producing one coherent dataset, and lead with the strongest real number
once you have it.}

\subsection*{Automated triage}
\tbd{Report the proportion of frames cleared without human attention. In the
synthetic demonstration set, 550 of 1,000 frames were empty and 450 animal
detections were distributed across 10 classes; these figures are illustrative
placeholders only.}

\subsection*{Validation effort}
\tbd{From a real run: images reviewed per hour in the gallery, or total reviewer
time for the dataset. This is the number that demonstrates the tool saves labour,
and it is the most persuasive single figure in the paper.}

\subsection*{Distributed validation}
\tbd{Report a real volunteer round if at all possible: number of volunteers,
number of batches, turnaround time, and the rate at which volunteers corrected
the model. Even a single small pilot round would be compelling evidence for the
central claim.}

\subsection*{Standards-compliant output}
\tbd{Confirm that a valid Camtrap DP package was produced, ideally by passing it
through the standard's validator, and report the outcome.} \cw

\subsection*{Robustness at scale}
\tbd{If available, report completion of the large real deployment, including at
least one checkpoint-and-resume event, as direct evidence that the field-scale
engineering works in practice.}

%% ---------------------------------------------------------------------------
\section{Discussion}
\note{Target 600-700 words.}

\subsection*{What is new}
Restate that the contribution is the end-to-end offline integration, and in
particular the built-in distributed volunteer validation loop, which no
comparable free desktop tool provides \cw.
\note{Claim nothing about model accuracy. That credit belongs to the developers of
the underlying models, and claiming it would be both wrong and easy to catch.}

\subsection*{Generalisability}
\note{Required by the article type: MEE asks that tools built for a specific
system be shown to adapt to more general use. Devote real space to this.}
The workflow is a template rather than a single-region application. The registry
allows any detector-classifier pair to be substituted, so a team working on
another continent, another taxon, or a locally fine-tuned model uses the same
pipeline unchanged \cw. Standards-compliant export means the result feeds
existing infrastructure regardless of region or taxon \cw. The volunteer loop is
content-agnostic, distributing whatever classes the chosen model produces.
\tbd{Give one or two concrete ``you could use this for X'' examples to make the
generality tangible rather than asserted.}

\subsection*{Limitations}
\note{State these plainly and first-hand. Reviewers will find them anyway, and
candour in a tools paper is persuasive rather than damaging.}
The application currently targets Windows only. The global model and the mapping
view require internet access, so only the European workflow is fully offline.
Validation quality depends on reviewer skill and, for volunteers, on task design
\cw. The bundled distribution is large. Most importantly, the tool orchestrates
rather than improves the underlying models, and therefore inherits their
taxonomic and geographic biases \cw.

\subsection*{Future directions}
Cross-platform builds; multi-observer agreement and consensus within the
volunteer loop \cw; optional active learning to prioritise uncertain images for
human review \cw; and additional models in the registry.

\subsection*{Conclusion}
\texttt{wildtag.ai} closes the last mile between an automated prediction and a
verified, shareable dataset, offline and without code, in a way that generalises
across models and regions.

%% ---------------------------------------------------------------------------
\section*{Author contributions}
\tbd{CRediT statement.}

\section*{Acknowledgements}
\tbd{Funders, volunteers, and the developers of the underlying models.}

\section*{Conflict of interest}
\tbd{None declared, or declare.}

\section*{Data availability statement}
\tbd{Repository and DOI for the code, the demonstration fixture and the
demonstration outputs.}

%% ---------------------------------------------------------------------------
\clearpage
\section*{Tables}

\begin{table}[htbp]
\centering
\caption{Capabilities of \texttt{wildtag.ai} compared with representative
camera-trap tools. \note{Verify every cell against the current release of each
tool before submission. The row that must survive scrutiny is the zero-install
distributed validation row, because it carries the paper's novelty claim.}}
\label{tab:compare}
\small
\begin{tabular}{p{5.1cm} c c c c c}
\toprule
Capability & wildtag.ai & Tool B & Tool C & Tool D & Tool E \\
\midrule
Runs offline, no account required        & Yes & Yes & No  & No  & Yes \\
No programming required                  & Yes & Yes & Yes & Yes & Yes \\
Multi-model registry (detector + classifier) & Yes & Yes & Fixed & Fixed & None \\
Structured in-application validation     & Yes & Partial & Yes & Yes & Yes \\
\textbf{Zero-install distributed volunteer validation} & \textbf{Yes} & No & No & Partial & No \\
Native standards-compliant export        & Yes & Partial & Export & Export & No \\
Checkpoint and resume for very large runs & Yes & \tbd{?} & n/a & n/a & n/a \\
Cost                                     & Free & Free & Free tier & Free/paid & Free \\
\bottomrule
\end{tabular}
\end{table}

\begin{table}[htbp]
\centering
\caption{Species summary for the demonstration deployment.
\note{Illustrative placeholder values from the synthetic fixture. Regenerate from
a real wildtag run on real images before submission.}}
\label{tab:species}
\small
\begin{tabular}{l r}
\toprule
Class & Detections \\
\midrule
Empty (no animal)  & \tbd{550} \\
Roe deer           & \tbd{104} \\
Red deer           & \tbd{77}  \\
Badger             & \tbd{54}  \\
Bird               & \tbd{50}  \\
Fox                & \tbd{46}  \\
Wild boar          & \tbd{37}  \\
Squirrel           & \tbd{27}  \\
Human              & \tbd{21}  \\
Lagomorph          & \tbd{18}  \\
\midrule
Total images       & \tbd{1{,}000} \\
Sites              & 10 \\
\bottomrule
\end{tabular}
\end{table}

%% ---------------------------------------------------------------------------
\clearpage
\section*{Figures}

\begin{figure}[htbp]
\centering
\fbox{\parbox[c][5cm][c]{0.85\textwidth}{\centering
\note{Figure 1 placeholder. Architecture and data flow: graphical interface,
subprocess inference engine, model registry, project folder, and outputs.}}}
\caption{Architecture and data flow of \texttt{wildtag.ai}.}
\label{fig:arch}
\end{figure}

\begin{figure}[htbp]
\centering
\fbox{\parbox[c][5cm][c]{0.85\textwidth}{\centering
\note{Figure 2 placeholder. The distribute-validate-collect-merge loop. This is
the paper's signature figure and the one reviewers will remember; invest real
effort in it.}}}
\caption{The distributed volunteer validation loop.}
\label{fig:loop}
\end{figure}

\begin{figure}[htbp]
\centering
\fbox{\parbox[c][5cm][c]{0.85\textwidth}{\centering
\note{Figure 3 placeholder. (a) The validation gallery with structured
correction. (b) The generated deployment map for the demonstration survey.}}}
\caption{(a) In-application validation gallery. (b) Generated deployment map.}
\label{fig:ui}
\end{figure}

%% ---------------------------------------------------------------------------
\clearpage

%% \nocite{*} forces the template entries in refs.bib to appear so that the
%% skeleton compiles with a visible reference list before any real citations
%% exist. DELETE this line as soon as you have real \citep/\citet citations in
%% the text, otherwise every entry in refs.bib will be printed whether cited
%% or not.
\nocite{*}

\bibliographystyle{besjournals}
\bibliography{refs}

\end{document}
